\subsection{Diffusion, scaling, and the determinant}\label{permanente}

In this subsection $K$ is the order of a doubly stochastic matrix $P$, and
$J_K=\mathbf1\mathbf1^T/K$. The estimates also apply with $K=k$ in Section~\ref{costes}.
The outer bounds give constants $M_0$ and $\rho_*<1$ such that
\begin{equation}\label{terminal-hipotesis-matriz}
 0\le P_{ij}\le M_0/K,\qquad \|P-J_K\|_{\rm op}\le\rho_* .
\end{equation}
Indeed, for a doubly stochastic matrix, the row and column block with weights $P$ is
$\left(\begin{smallmatrix}I&P\\P^T&I\end{smallmatrix}\right)$, with eigenvalues $1\pm\sigma_j(P)$.
The two lines of constant vectors correspond to the singular value one. A gap $\gamma>0$ off
the zero mode implies that the remaining singular values are at most $1-\gamma$, which is
precisely the required bound for $P-J_K$.

\begin{lema}[Doubly stochastic scaling]\label{terminal-escalado}
Suppose \eqref{terminal-hipotesis-matriz} holds, and write
$W_{ij}=P_{ij}(1-P_{ij})$ and $a_K=\sum_{ij}P_{ij}^2$. For all sufficiently large $K$, there are
vectors $\alpha,\beta$ with $\sum_i\alpha_i=\sum_j\beta_j$ such that the matrix with entries
$D_{ij}=e^{\alpha_i+\beta_j}W_{ij}$ is doubly stochastic. Moreover,
\[
 \max\{\|\alpha\|_\infty,\|\beta\|_\infty\}=O(K^{-1}),\qquad
 \sum_i\alpha_i+\sum_j\beta_j=a_K+O(K^{-1}),\qquad
 \|D-P\|_{\rm HS}=O(K^{-1}).
\]
The constants depend only on $M_0,\rho_*$. The support of $D$ coincides with that of $P$.
\end{lema}
\begin{proof}
To control the accuracy of the scaling, treated in general in \cite{sk}, consider $S=PP^T$.
It is symmetric and doubly stochastic, with $S_{ij}\le M_0/K$ and $\|S-J_K\|_{\rm op}\le\rho_*^2$,
since $S-J_K=(P-J_K)(P-J_K)^T$. Each row of $S-J_K$ has squared Euclidean norm
$\sum_j S_{ij}^2-1/K\le (M_0-1)/K$. Since $S^t-J_K=(S-J_K)^t$ for every integer $t\ge1$,
multiplying a row of $S-J_K$ by $(S-J_K)^{t-1}$ and passing from the Euclidean norm to the sum
of absolute values gives
\[
 \|S^t-J_K\|_{\infty\to\infty}
 \le \sqrt{M_0-1}\,\rho_*^{\,2(t-1)} .
\]
On mean-zero vectors, the series $\sum_{t\ge0}S^t$ is therefore the inverse of $I-S$, with
uniformly bounded $\infty\to\infty$ norm. The same proof with $P^TP$ gives the bound for
$(I-P^TP)^{-1}$ on this subspace.

Put $S_0=\left(\begin{smallmatrix}I&P\\P^T&I\end{smallmatrix}\right)$.
On the subspace $\sum x=\sum y$, consider the equation $S_0(x,y)=(f,g)$ for $\sum f=\sum g$.
If $s=K^{-1}\sum f$, elimination in the two blocks gives
\[
 \begin{aligned}
 x&=\tfrac12s\mathbf1+(I-PP^T)^{-1}(f-Pg),\\
 y&=\tfrac12s\mathbf1+(I-P^TP)^{-1}(g-P^Tf).
 \end{aligned}
\]
The arguments of both inverses have mean zero. The constant part comes from the fact that $x+Py=f$ and
$y+P^Tx=g$ both have mean $s$ and the means of $x,y$ agree. Since
$\|P\|_{\infty\to\infty}=\|P^T\|_{\infty\to\infty}=1$, the inverse $T$ of $S_0$ in this gauge
satisfies $\|T\|_{\infty\to\infty}\le C_1(M_0,\rho_*)$.

Let $\Psi(\alpha,\beta)$ be the vector of row and column sums of
$(e^{\alpha_i+\beta_j}W_{ij})$ minus one in each coordinate. The sums of its two blocks are
equal. At zero the defect is at most $M_0/K$ by $\sum_jP_{ij}^2\le M_0/K$ and its column
analogue. On a ball of radius $r\le1$, the diagonal blocks of the derivative of $\Psi$
contain the row and column sums of the matrix scaled at the point in question. The off-diagonal
blocks are that matrix and its transpose. In each row and column we have
\[
 \sum_j\big|e^{\alpha_i+\beta_j}W_{ij}-P_{ij}\big|
 \le C_2(r+K^{-1})\sum_jP_{ij}=C_2(r+K^{-1}),
\]
with the same inequality on summing over $i$. Here we used $|\alpha_i+\beta_j|\le2r$ and
$P_{ij}\le M_0/K$. Thus $\|\Psi'-S_0\|_{\infty\to\infty}\le 2C_2(r+K^{-1})$.

With the gauge fixed, the map $u\mapsto u-T\Psi(u)$ moves the origin by at most $C_1M_0/K$.
Take the closed ball of radius $r=2C_1M_0/K$. For large $K$ the derivative of the map has norm
at most $2C_1C_2(r+K^{-1})\le1/2$. The map therefore preserves this ball and is a contraction.
Its fixed point solves $\Psi=0$ and has norm $O(K^{-1})$. Since $P_{ij}<1$ for large $K$, the
factor $(1-P_{ij})e^{\alpha_i+\beta_j}$ is positive and preserves the support.

Uniformly in the indices, we now have
\[
 D_{ij}=P_{ij}\bigl(1+\alpha_i+\beta_j-P_{ij}+O(K^{-2})\bigr).
\]
Sum over $i,j$. Since both $D$ and $P$ have total mass $K$ and the rows and columns of $P$ sum
to one, we obtain $\sum_i\alpha_i+\sum_j\beta_j-a_K=O(K^{-1})$. Also,
$|D_{ij}-P_{ij}|\le C_3P_{ij}/K$ and $a_K\le M_0$. Squaring and summing proves the claim in
Hilbert--Schmidt norm.
\end{proof}
